% !TeX root = main.tex
\chapter*{ВСТУП}
\addcontentsline{toc}{chapter}{ВСТУП}

\textbf{Обґрунтування вибору теми дослідження.}
Практична проблема полягає в тому, що мультимедійні дані надходять до вебсистем через штатні канали завантаження, зберігання й поширення, але на подальшому маршруті їх інтерпретують компоненти з різними правилами визначення типу, меж контейнера й допустимих операцій. Після приймання той самий носій перевіряють, декодують, масштабують, повторно кодують, доповнюють службовими даними та передають між прикладними компонентами. Через це формальна коректність на вході характеризує лише придатність даних до певної операції й не виключає прихованого впливу на інший засіб оброблення.

У роботі розмежовано чотири близькі, але нетотожні поняття:

\begin{enumerate}[label=\arabic*),leftmargin=1.1cm]
\item Аномалія --- спостережуване відхилення від установленої структурної, статистичної або поведінкової норми. Окрема аномалія не доводить наявності загрози.
\item Аномальний код --- цілеспрямовано сформована послідовність даних, команда, фрагмент коду або структурне значення, що не відповідає заявленому призначенню носія та може підтримати несанкціоновану дію.
\item Прихована загроза --- сценарій можливого небезпечного впливу, для якого встановлюється зв’язок між властивостями носія, операцією оброблення та очікуваним наслідком.
\item Шкідливий код --- код, програмну семантику й шкідливу функцію якого підтверджено.
\end{enumerate}

Таке розмежування запобігає ототожненню статистичного відхилення з доведеним виконанням програми.

Прихований вплив може залишати різні спостережувані прояви. Просторове або частотне вбудовування змінює декодоване зображення; структурне розміщення використовує службові поля, межі елементів чи додаткові області контейнера; багатоформатний об’єкт допускає різну інтерпретацію одних байтів; поведінковий прояв фіксується у подіях лише під час декодування, перетворення або подальшого використання. Для відео до цього додаються часовий розподіл змін і залежність від структури стисненого потоку. Отже, зорова правдоподібність, коректний заявлений тип і успішне звичайне відтворення контролюють лише частину можливих проявів.

\textbf{Наукова проблема та дослідницька прогалина.}
Межа відомих засобів визначається представленням, яке вони спостерігають. Стегоаналітичні методи виявляють слабкі статистичні зміни, але не обов’язково аналізують службові області контейнера. Структурні правила фіксують порушення формату, проте не розпізнають адаптивне вбудовування у природні ділянки зображення. Аналіз подій і поведінки може зареєструвати нетипову подію, однак без причинної прив’язки до конкретного носія не пояснює її походження. Оцінки цих каналів також не можна механічно порівнювати, оскільки вони залежать від формату, джерела початкових об’єктів, способу приховування, складу набору даних і правила прийняття рішення.

Наукова проблема полягає у формуванні обґрунтованого рішення щодо стану мультимедійного об’єкта за неповних, різнорідних і залежних від маршруту спостережень. За результатами аналізу, поданого в розділі~\ref{chap:analysis}, дослідницьку прогалину визначено як недостатню опрацьованість узгодженого представлення, яке пов’язує структурні свідчення, результати спектрального й зорового аналізу та відомості про події й поведінку з маршрутом конкретного об’єкта у вебсистемі.

Наукова суперечність полягає в багаторівневому прояві прихованої загрози за переважно ізольованого формування локальних оцінок. Використання лише найбільшої локальної оцінки робить рішення чутливим до випадкового сильного відхилення. Вимога одночасного підтвердження всіма каналами, навпаки, створює пропуск за недоступної модальності або неактивованого прихованого механізму. Тому узгодження має зберігати сильне свідчення окремого джерела, позначати відсутні спостереження та відокремлювати оцінку ризику від упевненості в її підставах.

Дослідницьке припущення полягає в тому, що узгодження ознак за спільним ідентифікатором, походженням, локалізацією та маскою доступності дає змогу формувати простежуване рішення без прирівнювання відсутньої модальності до нормального стану. Для реалізації цього припущення потрібні модель, два взаємопов’язані методи й узгоджений інформаційний процес. Мультимедійний об’єкт до завершення перевірки має залишатися недовіреним; незмінний початковий об’єкт, декодоване представлення, структурна карта й причинно пов’язані події не повинні змішуватися. Статичні свідчення та відомості про події й поведінку необхідно інтегрувати до прийняття рішення, а технологічну дію визначати після оцінювання ризику й упевненості. Реєстрація результату повинна зберігати підстави рішення, версії компонентів і зв’язок із початковим об’єктом.

Перевірюваними наслідками припущення є збереження внеску й доступності кожної групи ознак, відтворюваність переходу від спостережень до аналітичного стану та можливість окремо оцінити помилки рішення. Агреговані результати попередньої експериментальної серії дають змогу ретроспективно перевірити лише описану прототипну конфігурацію. Детерміновані структурні, спектральні й зорові ознаки першого методу оцінено у внутрішньому групово відокремленому експерименті, а перенесення незмінних параметрів --- на чотирьох зовнішніх наборах, що разом містили $38$ реальних відеозаписів із $12$ пристроїв та $24$ груп походження. Скорочена модель узгодження статичних оцінок із урахуванням формату опрацьовувала лише вже сформовані оцінки першого методу і тому не є перевіркою повного методу поведінкового узгодження. Широкої переносимості на довільні джерела, кодеки й профілі форматів, а також кількісної переваги повного мультимодального контуру не встановлено.

Потреба в дослідженні зумовлена поєднанням поширеного мультимедійного обміну, багаторівневої будови носія та різних правил його інтерпретації компонентами вебсистеми. За цих умов контроль одного представлення залишає неперевіреними інші канали, що обмежує обґрунтованість рішення щодо стану об’єкта. Формування інформаційної технології для узгодження таких свідчень належить до задач спеціальності 126 «Інформаційні системи та технології», оскільки охоплює моделювання інформаційного процесу, розроблення методів аналізу, організацію технологічного контуру й експериментальне оцінювання його складників.

\textbf{Зв’язок роботи з науковими програмами, планами і темами.}
Дисертаційну роботу виконано в Хмельницькому національному університеті під час підготовки за спеціальністю 126 «Інформаційні системи та технології». Формальний зв’язок із зареєстрованою науково-дослідною темою в роботі не заявляється, оскільки в наявних матеріалах немає її назви, номера державної реєстрації, строку виконання та документально підтвердженої ролі здобувача.

\textbf{Мета і завдання дослідження.}
Метою дослідження є формування інформаційної технології виявлення визначених структурних, спектральних і зорових ознак, а також ознак подій і поведінки, що характеризують аномальний вміст у мультимедійних даних вебсередовища, на основі моделі мультимодального представлення прихованої загрози, методу аналізу спектральних, зорових і структурних ознак та методу поведінкового узгодження ознак різного походження. Експериментальна частина перевіряє виявлення контрольовано внесених структурних і статистичних змін та не ототожнює позитивне рішення з доведеною наявністю виконуваного або шкідливого коду.

Для досягнення поставленої мети необхідно розв’язати такі завдання:

\begin{enumerate}[label=\arabic*)]
\item проаналізувати використання мультимедійних об’єктів у вебсистемах як засобів прихованого впливу, способи прихованого розміщення даних і відомі методи виявлення відповідних ознак та визначити взаємодоповнювальні обмеження цих методів;

\item формалізувати процес використання мультимедійних даних для передавання прихованого впливу через штатні операції вебсистеми;

\item розробити модель мультимодального представлення прихованої загрози, яка інтегрує структурні, спектральні й зорові ознаки, а також ознаки подій і поведінки на основі багаторівневого інформаційного контейнера, маски доступності, критеріїв рішення та процесу виявлення;

\item удосконалити метод аналізу спектральних, зорових і структурних ознак мультимедійних даних;

\item розробити метод поведінкового узгодження ознак різного походження для зіставлення статичних ознак із причинно пов’язаними подіями оброблення;

\item сформувати п’ятиетапну інформаційну технологію виявлення ознак аномального вмісту в мультимедійних даних вебсередовища;

\item обґрунтувати структуру прототипної конфігурації, зафіксованої в контрольному описі експериментального запуску від 11 липня 2026~р., та визначити склад експериментального запису, потрібного для повторної перевірки;

\item провести експериментальне оцінювання детермінованої статичної частини першого методу та скороченої моделі узгодження статичних оцінок у шести відокремлених серіях --- ретроспективній, внутрішній групово відокремленій, першій зовнішній, поетапній підтверджувальній, великомасштабній багатоформатній і заздалегідь відокремленій зовнішній, --- доповнити показники рішень окремим ресурсним випробуванням і локальним аудитом станів $S_1$--$S_5$, визначити співвідношення достовірності виявлення й обчислювальних витрат та встановити межі допустимого тлумачення одержаних результатів.
\end{enumerate}

Простежуваність поставлених завдань наведено в табл.~\ref{tab:intro_traceability}. Позначення доказового статусу відокремлює формальний результат від результату, підтвердженого експериментом.

\begin{longtable}{@{}>{\centering\arraybackslash}p{0.07\textwidth}p{0.27\textwidth}p{0.28\textwidth}p{0.29\textwidth}@{}}
\caption{Відповідність завдань, результатів і доказового статусу}\label{tab:intro_traceability}\\
\toprule
№ & Результат і місце викладу & Експериментальна опора & Доказовий статус \\
\midrule
\endfirsthead
\multicolumn{4}{l}{Продовження таблиці \thetable}\\
\toprule
№ & Результат і місце викладу & Експериментальна опора & Доказовий статус \\
\midrule
\endhead
\bottomrule
\endfoot
1 & Критичний аналіз способів прихованого впливу і засобів виявлення, розділ~\ref{chap:analysis} & Систематизація відомих праць & Аналітично виконано \\
2 & Причинний маршрут носія та умови прихованого впливу, підрозділ~\ref{sec:attack_model} & Перевірка логічної узгодженості сценаріїв & Формалізовано \\
3 & Модель представлень, доступності, ризику й рішення, розділ~\ref{chap:models_formalization} & Непряма перевірка складників у розділі~\ref{chap:experiments} & Формалізовано; повний контур не перевірено \\
4 & Метод статичного аналізу, підрозділ~\ref{sec:svs_method} & Внутрішня та чотири зовнішні серії & Перевірено для заявлених змін і п’яти форматів \\
5 & Метод поведінкового узгодження ознак різного походження, підрозділ~\ref{sec:behavior_alignment} & Лише скорочена ознака, сформована за правилами; повних журналів немає & Розроблено формальний опис; повну конфігурацію не перевірено \\
6 & П’ятиетапна технологія, підрозділ~\ref{sec:information_technology} & Ресурсне випробування та локальний аудит $45$ сценарних спостережень зі штатними, відмовними, повторними й паралельними циклами & Локальні переходи й відмови перевірено; промисловий контур не підтверджено \\
7 & Специфікація конфігурації та експериментального запису, розділ~\ref{chap:experiments} & Машинно перевірені реєстри, ознаки, прогнози, параметри та контрольні суми нових серій; ретроспективний пакет і первинний модуль виділення спектральних ознак неповні & Числові підсумки нових серій відтворено; наскрізне історичне виділення ознак відтворено частково \\
8 & Шість серій, ресурсне випробування й аудит станів, розділ~\ref{chap:experiments} & Груповий поділ, зовнішні набори, пооб’єктні рішення нових серій і машинозчитувані записи $S_1$--$S_5$ & Експериментально виконано в установлених межах \\
\end{longtable}

\textbf{Об’єкт дослідження.}
Об’єктом дослідження є процес формування рішення щодо ознак аномального вмісту під час приймання й оброблення мультимедійних даних компонентами вебсистеми.

\textbf{Предмет дослідження.}
Предметом дослідження є модель мультимодального представлення прихованої загрози, метод аналізу спектральних, зорових і структурних ознак, метод поведінкового узгодження ознак різного походження та інформаційна технологія виявлення ознак аномального вмісту в мультимедійних даних вебсередовища.

\textbf{Методи дослідження.}
Системний аналіз використано для встановлення меж досліджуваного процесу та причинних зв’язків між носієм, операцією оброблення і спостережуваним результатом. Методи математичного моделювання застосовано для формалізації маршруту прихованого впливу, багаторівневого контейнера, мультимодального представлення, оцінки ризику, упевненості й правила прийняття рішення. Результатом їх застосування є модельний контракт, який розділяє вхідні представлення, аналітичний стан і технологічну дію.

Для формування статичного опису використано методи цифрового оброблення зображень і відео та структурного аналізу контейнерів. Перші утворюють високочастотні залишки, характеристики дискретного косинусного й дискретного хвилькового перетворень, просторові локалізації та часові агрегати кадрів. Другі перевіряють заголовки, межі елементів, службові поля, завершальні маркери й додаткові байтові області. Детерміновані ознаки перевірено внутрішнім групово відокремленим протоколом для PNG і MP4, а перенесення незмінних параметрів першого методу досліджено на реальних відеозаписах у форматах PNG, JPEG, WebP, MP4 і WebM. Методи машинного навчання застосовано для контрольних моделей статичних ознак та моделі узгодження вихідних оцінок першого методу з урахуванням формату. Повну модель зорового перетворювача не навчено й незалежно не перевірено.

Для пов’язування статичного опису з реакцією середовища використано причинний аналіз подій і методи мультимодального узгодження за спільним ідентифікатором, часовими межами та маскою доступності. В експериментальних серіях доступна технічна ознака не спрацювала у $4\,832$ унікальних відеоматеріалах, а причинно пов’язаних журналів не сформовано. Тому скорочену модель узгодження статичних оцінок із урахуванням формату не можна вважати реалізацією повного методу поведінкового узгодження. Кодувальник послідовності подій і перехресне зіставлення формалізовано, але не навчено й не перевірено незалежно.

Методи арифметичної перевірки та описового статистичного оцінювання застосовано до збереженої агрегованої матриці помилок, компонентних конфігурацій і часових характеристик описаного прототипу. Для внутрішньої серії використано груповий поділ, групове повторне вибирання та парне зіставлення конфігурацій. У чотирьох зовнішніх наборах незмінні параметри переносили на нові записи, а залежність похідних об’єктів враховували за групами походження. Фактичні ресурсні витрати оцінено за 95-м процентилем пооб’єктного наскрізного часу та найбільшим обсягом фізичної пам’яті процесу у п’яти окремих однопотокових запусках. Розвідувальний інтегральний рейтинг $K_I$ об’єднує чотири складники: достовірність виявлення, нормовані часові витрати, нормовані витрати пам’яті та форматну стійкість. Завершеність, цілісність і мінімальну форматну придатність перевіряли окремо від $K_I$, щоб висока швидкодія не приховала неприйнятного результату окремого формату. Ретроспективна процедура дає змогу відтворити доступні похідні показники, але не відновлює відсутніх пооб’єктних рішень або повних журналів попереднього запуску.

\textbf{Наукова новизна одержаних результатів.}
Наукову новизну становлять чотири взаємопов’язані результати:

\begin{enumerate}[label=\arabic*)]
\item \textbf{розроблено модель мультимодального представлення прихованої загрози}, яка подає мультимедійний об’єкт як багаторівневий інформаційний контейнер. На відміну від підходів, що зіставляють локальні оцінки без збереження їхнього походження, модель інтегрує структурні, спектральні й зорові ознаки, а також ознаки подій і поведінки через спільний ідентифікатор, маску доступності та відомості про джерело. Це дає змогу формувати типізоване представлення $Z$, зберігати локалізації й посилання на причинно пов’язаний журнал та відокремлювати аналітичний висновок від технологічної дії;

\item \textbf{удосконалено метод аналізу спектральних, зорових і структурних ознак мультимедійних даних}, який узгоджує детерміновані залишкові, частотні, хвилькові та часові характеристики з результатами структурного розбору контейнера. На відміну від окремого аналізу декодованого подання або контейнера, метод зберігає оцінки й локалізації двох статичних гілок до їх інтеграції, ураховує доступність ознак і застосовує форматні правила без перенесення неперевірених властивостей між форматами. Науковий акцент зроблено на двогілковій структурі зі збереженням походження свідчень, а не на спектральному каналі як самодостатньому засобі виявлення: зовнішня перевірка на реальних відеозаписах показала межу ізольованих спектральних і зорових ознак, тоді як структурна гілка та спільна числова проєкція $Z_c$ зберегли діагностичну цінність у встановлених межах. Детерміновані складники первинно оцінено для PNG і MP4, а перенесення незмінних параметрів на JPEG, WebP і WebM перевірено без повторного навчання; універсальність для довільних кодеків, профілів і реалізацій форматів не заявляється;

\item \textbf{розроблено метод поведінкового узгодження ознак різного походження}, який задає зіставлення статичного опису з подіями контрольованого оброблення того самого об’єкта за ідентифікатором, часовими межами та фазовою належністю. На відміну від загального послідовного застосування статичного й динамічного аналізу, метод установлює вхідний і вихідний контракти, умови причинного допуску події, граничні стани, зберігає посилання на початковий журнал і відрізняє спостережувану відсутність подій від недоступності журналу. Цей результат має формальний статус: повну конфігурацію з кодувальником послідовностей і перехресним зіставленням не реалізовано й експериментально не підтверджено;

\item \textbf{набула подальшого розвитку інформаційна технологія виявлення ознак аномального вмісту в мультимедійних даних вебсередовища}, у якій модель і два методи організовано в п’ять послідовних етапів: попереднє оброблення, формування ознак, інтеграція ознак, прийняття рішення та реєстрація результатів. Відмінність технології полягає не в числі етапів, ізоляції чи реконструкції як таких, а в типізованих переходах між станами, незмінному збереженні початкового об’єкта, відокремленні аналітичного рішення від дії політики, реконструкції як нового циклу та реєстраційному записі до виконання зовнішньої дії. Локально перевірено переходи $S_1$--$S_5$, атомарний запис без перезапису, повтор, відновлення та безпечне завершення за двох класів технічних відмов; повний поведінковий і промисловий контури не підтверджено.
\end{enumerate}

\textbf{Практичне значення одержаних результатів.}
Практичне значення полягає у визначенні структури технологічного контуру між штатним каналом приймання мультимедійного об’єкта та прикладною логікою вебсистеми. Контур передбачає ізоляцію незмінного входу, структурний аналіз, аналіз спектральних і зорових ознак, контрольоване спостереження, інтеграцію свідчень, формування рішення та контрольну перевірку. Типізований вихід призначено для вибору однієї з обґрунтованих політикою дій: допуску, експертної перевірки, реконструкції нового об’єкта або блокування.

У контрольному описі експериментального запуску зафіксовано дослідний прототип статичної конфігурації та скороченої подієвої оцінки, сформованої за правилами. Позитивна еталонна мітка в контрольному наборі позначала внесене структурне або статистичне відхилення, а не доведену наявність виконуваного аномального коду. Для $110$ об’єктів з агрегованої матриці помилок арифметично одержано частку хибнопозитивних рішень $0{,}1273$ і частку хибнонегативних рішень $0{,}0909$.

У внутрішньому відтворюваному експерименті для першого методу сформовано $200$ первинних груп PNG і MP4 та $400$ похідних об’єктів із груповим поділом $60:20:20$. На відокремленій тестовій частині з $80$ об’єктів, яку не використовували під час навчання й налаштування, спільна числова проєкція $Z_c$ дала $40$ правильно визначених контрольних, жодного хибнопозитивного, два хибнонегативні та $38$ правильно визначених змінених об’єктів; частка правильних рішень становила $0{,}9750$, повнота --- $0{,}9500$, $F_1=0{,}9744$, а площа під характеристичною кривою --- $0{,}9869$.

Для зіставлення достовірності виявлення й ресурсних обмежень обчислено розвідувальний інтегральний рейтинг $K_I$. Він поєднує середню за джерелами збалансовану частку правильних рішень, нормовані 95-й процентиль часу та найбільший обсяг фізичної пам’яті процесу, а також стійкість між форматами. Для структурної конфігурації, конфігурації спектральних і зорових ознак, спільної та узгоджувальної конфігурацій одержано відповідно $77{,}50$, $50{,}69$, $68{,}95$ і $69{,}46$ бала. Окрему форматну умову виконали лише спільна й узгоджувальна конфігурації; тому $K_I$ використано разом із первинними складниками, а не як самодостатню міру якості технології.

Чотири зовнішні набори разом охопили $38$ реальних відеозаписів із $12$ пристроїв і $24$ груп походження. Перша серія на $56$ похідних MP4 не відтворила переваги спільної проєкції над структурною конфігурацією: збалансована частка правильних рішень становила відповідно $0{,}7857$ і $0{,}8333$, тоді як для конфігурації спектральних і зорових ознак --- $0{,}5000$. Поетапна серія на $32$ об’єктах установила, що правило підтвердження усуває чотири хибнопозитивні рішення, але збільшує кількість пропусків від чотирьох до восьми. У багатоформатній серії $500$ базових фрагментів породили $10\,000$ об’єктів форматів PNG, JPEG, WebP, MP4 і WebM; це число не прирівнювалося до кількості незалежних джерел. У заздалегідь відокремленій зовнішній серії на $2\,000$ об’єктах із чотирьох нових пристроїв точкова оцінка збалансованої частки правильних рішень змінилася від $0{,}7827$ для спільної конфігурації до $0{,}8307$ для моделі узгодження. Вона усунула $76$ хибнопозитивних рішень, але додала $84$ пропуски та не виявила жодної з $500$ суто спектральних змін; критерій Мак-Немара дав $p=0{,}5801$.

Окремий аудит повторного виділення ознак із незмінених $56$ і $2\,000$ MP4 повністю зберіг структурні рішення, але виявив калібрувальну нестійкість спектральної та спільної проєкцій. Тому архівні значення пов’язано з канонічними таблицями ознак, а поточну еталонну реалізацію не подано як чисельно тотожну історичній. Парний розвідувальний аналіз $2\,500$ форматних представлень установив ненульову декодовану спектральну відмінність у кожній парі та піксельну тотожність усіх $5\,000$ контейнерних контрольних пар. Отже, повне зникнення сліду не є єдиним поясненням низької якості спектральної гілки, однак збереження первинної просторово-часової локалізації не встановлено.

Окремий розвідувальний аналіз $1\,500$ допустимо змінених контрольних об’єктів установив недостатню специфічність структурної гілки. Для приватних розширень або відступів частка хибнопозитивних рішень дорівнювала $1{,}0000$ у кожному з п’яти форматів, хоча всі об’єкти пройшли розбір і декодування, а $1\,000$ контейнерних пар були піксельно тотожними. Отже, структурна оцінка є формальним індикатором, а не самодостатнім доказом прихованого впливу. Її тлумачення потребує перевірки відповідності формату й правилам безпеки, а також експертної оцінки.

Локальний аудит інформаційної технології охопив $45$ сценарних спостережень для десяти незмінних об’єктів п’яти форматів. У штатних, повторних, відновлених і чотирипроцесних циклах створено $35$ цілісних записів $S_5$; п’ять помилок декодування і п’ять контрольовано внесених відмов реєстрації завершилися без $S_5$ та без виконання зовнішнього припису. Медіана повного послідовного циклу становила $20{,}321$~мс, 95-й процентиль --- $108{,}949$~мс, а один чотирипроцесний запуск дав $22{,}827$ об’єкта за секунду за найбільшої сукупної пам’яті $545{,}52$~МіБ. Ці показники підтверджують локальну реалізацію переходів і відмов, але не розподілену або промислову масштабованість.

Показники ретроспективної серії характеризують лише зафіксовану конфігурацію на наборі, використаному також для налаштування параметрів. Внутрішній групово відокремлений цикл усунув змішування похідних спільного походження, а зовнішні набори частково перевірили перенесення незмінних параметрів на реальні відеозаписи й п’ять форматів. Водночас тисячі форматно-сценарних похідних походять від обмеженої кількості базових фрагментів і пристроїв, тому широкої переносимості не встановлено. Доступна технічна поведінкова ознака не спрацювала у $4\,832$ унікальних відеоматеріалах, а причинно пов’язаних журналів $X_b$ не сформовано; повний другий метод залишається неперевіреним. Прототип не є завершеним виробничим продуктом, тож обґрунтоване застосування обмежується ізоляцією та експертною перевіркою, а безумовне автоматичне блокування наявними даними не підтримується.

\textbf{Особистий внесок здобувача.}
Дисертація узагальнює результати дослідження здобувача у вигляді однієї моделі, двох методів та інформаційної технології. Положення, використані з праць інших авторів, супроводжено відповідними посиланнями; спільні публікації наведено у списку праць здобувача. Відповідність положень новизни публікаціям і наявним експериментальним матеріалам подано в додатку~\ref{app:publications_and_approbation}. Конкретний розподіл постановки задачі, формалізації, програмної реалізації, проведення експерименту, статистичного опрацювання та підготовки тексту між співавторами не заявляється без авторських довідок; ці відомості мають бути внесені до остаточної редакції після документального погодження.

\textbf{Апробація результатів дисертації.}
Матеріали, що засвідчують апробацію результатів, опубліковано у виданнях таких наукових заходів: 13-ї IEEE International Conference on Dependable Systems, Services and Technologies (DESSERT 2023; Афіни, 13--15 жовтня 2023~р.); IntelITSIS 2024 (Хмельницький, 28 березня 2024~р.); ICyberPhyS 2024 (Хмельницький, 28 червня 2024~р.); 14-ї IEEE International Conference on Dependable Systems, Services and Technologies (DESSERT 2024; Афіни, 11--13 жовтня 2024~р.); AdvAIT 2024 (Хмельницький / Жиліна, 5 грудня 2024~р.); 13-ї Міжнародної наукової конференції «Безпека інформаційних технологій» (ITSec~2024; Львів, 9--11 травня 2024~р.); IntelITSIS 2025 (Хмельницький, 4 квітня 2025~р.).

\textbf{Публікації.}
За темою дисертації опубліковано дев’ять наукових праць. Дві статті містять основні наукові результати й опубліковані у фахових наукових виданнях України. Сім праць опубліковано у виданнях IEEE, CEUR Workshop Proceedings і матеріалах конференції «Безпека інформаційних технологій» (ITSec). Повні бібліографічні описи наведено у списку публікацій здобувача.

\textbf{Структура та обсяг дисертації.}
Дисертація складається з анотацій українською та англійською мовами, списку публікацій здобувача, змісту, переліку умовних позначень, вступу, чотирьох розділів, загальних висновків, списку використаних джерел і додатка. У першому розділі проаналізовано способи прихованого впливу й методи виявлення. У другому розділі сформовано модель мультимодального представлення прихованої загрози. У третьому розділі удосконалено метод аналізу спектральних, зорових і структурних ознак, розроблено метод поведінкового узгодження ознак різного походження, а п’ятиетапна інформаційна технологія набула подальшого розвитку.

У четвертому розділі проведено шість відокремлених експериментальних серій, окреме ресурсне випробування та локальний аудит станів $S_1$--$S_5$: ретроспективну перевірку $110$ об’єктів, внутрішню групово відокремлену перевірку $400$ PNG- та MP4-об’єктів і чотири зовнішні серії на $38$ реальних відеозаписах із $12$ пристроїв. Установлено межі форматного перенесення статичних конфігурацій, виконано заздалегідь відокремлену зовнішню перевірку скороченої моделі узгодження статичних оцінок, проведено аудити чутливості повторного виділення ознак, специфічності структурної гілки, збереження спектрального сліду й локальних технологічних відмов та кількісно описано співвідношення достовірності виявлення й обчислювальних витрат. Актуальні кількісні характеристики обсягу рукопису визначаються за остаточною зібраною редакцією.

\clearpage
